\chapter{Chest Pressure}

\section*{Definition}
Chest pressure is a sensation of squeezing, heaviness, tightness, fullness, or weight in the chest. It may arise from the heart, lungs, esophagus, muscles, nerves, or anxiety. New, severe, exertional, or persistent chest pressure must be treated as potentially cardiac until evaluated.

\section*{What Happens in Your Body}
Heart ischemia can produce pressure when oxygen demand exceeds coronary blood supply. Lung disease can increase the work of breathing, while reflux, esophageal spasm, chest-wall inflammation, and muscle strain stimulate nearby pain pathways. Cardiac ischemia may be atypical, particularly in older adults, women, and people with diabetes.

\section*{Causes}
\begin{itemize}
  \item Acute coronary syndrome, stable angina, coronary vasospasm, or microvascular angina
  \item Pulmonary embolism, pneumothorax, pneumonia, asthma, or pulmonary hypertension
  \item Aortic dissection, pericarditis, or myocarditis
  \item Gastroesophageal reflux, esophageal spasm, or peptic disease
  \item Costochondritis, muscle strain, rib injury, panic attack, or anxiety
\end{itemize}

\section*{When to See a Doctor}
Call emergency services for pressure lasting more than a few minutes, recurring, or accompanied by shortness of breath, sweating, nausea, fainting, palpitations, or pain radiating to the arm, jaw, back, or upper abdomen. Do not drive yourself and do not wait to see whether an antacid helps.

\section*{Related Symptoms}
\begin{itemize}
  \item Breathlessness, fatigue, dizziness, palpitations, nausea, or sweating
  \item Cough, wheezing, fever, pain with breathing, or coughing blood
  \item Tenderness with movement or touch, heartburn, or sour regurgitation
\end{itemize}

\section*{Self-Care / Home Management}
Stop exertion, sit upright, and call emergency services for concerning symptoms. Use prescribed rescue medicine only according to the existing action plan. Do not take someone else’s cardiac medication or delay assessment with home remedies.

\section*{How Your Doctor Will Evaluate You}
\begin{itemize}
  \item Immediate assessment includes vital signs, oxygen saturation, cardiovascular examination, and a 12-lead ECG.
  \item Serial high-sensitivity troponin testing evaluates myocardial injury; blood count, electrolytes, glucose, and chest imaging may be added.
  \item CT angiography, echocardiography, stress testing, or coronary angiography is selected according to risk and initial findings.
\end{itemize}

\section*{Treatment Options}
Treatment depends on the cause and may include emergency coronary reperfusion, antiplatelet or anticoagulant therapy, treatment of pulmonary embolism, antibiotics, bronchodilators, acid suppression, or musculoskeletal care. Cardiac causes require clinician-directed risk-factor management and follow-up.

\section*{References}
\begin{thebibliography}{99}
\bibitem{chest_pressure:ref1} Gulati M, et al. 2021 AHA/ACC Guideline for the Evaluation and Diagnosis of Chest Pain. \textit{Circulation}. 2021;144:e368--e454.
\bibitem{chest_pressure:ref2} Byrne RA, et al. 2023 ESC Guidelines for the management of acute coronary syndromes. \textit{Eur Heart J}. 2023;44:3720--3826.
\bibitem{chest_pressure:ref3} Konstantinides SV, et al. 2019 ESC Guidelines for acute pulmonary embolism. \textit{Eur Heart J}. 2020;41:543--603.
\end{thebibliography}
